\documentclass[11pt]{article}
\usepackage[letterpaper,margin=0.72in]{geometry}
\usepackage{amsmath,amssymb,amsthm,mathtools}
\usepackage{graphicx}
\usepackage[T1]{fontenc}
\usepackage{microtype}
\setlength{\parindent}{0pt}
\setlength{\parskip}{0.42em}
\setlength{\emergencystretch}{2em}
\newtheorem{theorem}{Theorem}
\newtheorem{corollary}{Corollary}
\newcommand{\Lip}{\operatorname{Lip}}
\newcommand{\R}{\mathbb R}

\title{The pointed Lipschitz space need not be complete\\
\large A negative answer to Remark 3.3 of arXiv:2410.10677}
\author{Run \texttt{fa\_banach\_001} --- \texttt{agent\_lane\_03}}
\date{13 August 2026}

\begin{document}
\maketitle

\textbf{Status.} High-confidence candidate full counterexample, pending
human review.  The proposed completeness fails already for maps
$\R\to\R$.

\section*{The source question}

Let $M,N$ be metric spaces with distinguished point $0$.  Write
$\Lip_0(M,N)$ for the Lipschitz maps $f:M\to N$ with $f(0)=0$, and
\[
 d_e(f,g)=\sup_{x\ne0}\frac{d(f(x),g(x))}{d(x,0)}.
\]
Karn and Mandal prove that the larger space $E(M,N)$ of extensively bounded
maps is complete whenever $N$ is complete.  Their Remark 3.3 asks whether
$\Lip_0(M,N)$ is also complete under the restricted metric $d_e$.

\begin{center}
\includegraphics[width=0.84\textwidth]{figures/open_problem_crop.png}
\end{center}
The exact source page is bundled as
\texttt{source\_excerpt\_remark\_3\_3\_page\_11.pdf}.

\begin{theorem}\label{thm:main}
With the usual metric and distinguished point zero,
$\Lip_0(\R,\R)$ is not complete under $d_e$.  More precisely, there are
compactly supported Lipschitz functions $f_n:\R\to\R$ that are $d_e$-Cauchy
and converge in $E(\R,\R)$ to a continuous extensively bounded function
that is not Lipschitz.
\end{theorem}

\section*{Construction and proof}

Put $a=\tfrac12$ and
\[
 \eta(x)=\bigl(1-4|x-a|\bigr)_+,
 \qquad
 r_n(x)=\bigl((x-a)^2+n^{-2}\bigr)^{1/4},
 \qquad
 r(x)=|x-a|^{1/2}.
\]
Define
\[
 f_n(x)=\eta(x)r_n(x),\qquad f(x)=\eta(x)r(x).
\]
The tent function $\eta$ is $4$-Lipschitz and is supported in
$[\tfrac14,\tfrac34]$.  For each fixed $n$, $r_n$ is continuously
differentiable and
\[
 |r_n'(x)|
 =\frac{|x-a|}{2((x-a)^2+n^{-2})^{3/4}}
 \le \frac{\sqrt n}{2}.
\]
Thus $f_n$ is a compactly supported Lipschitz function.  Since its support
misses zero, $f_n(0)=0$, so $f_n\in\Lip_0(\R,\R)$.

For $u,v\ge0$ and $0<\alpha<1$, subadditivity gives
$(u+v)^\alpha\le u^\alpha+v^\alpha$.  Taking
$u=(x-a)^2$, $v=n^{-2}$, and $\alpha=\tfrac14$ yields
\[
 0\le r_n(x)-r(x)\le n^{-1/2}.
\]
Outside $[\tfrac14,\tfrac34]$ both $f_n$ and $f$ vanish, while on this
interval $|x|\ge\tfrac14$.  Consequently
\begin{equation}\label{conv}
 d_e(f_n,f)
 =\sup_{x\ne0}\frac{|f_n(x)-f(x)|}{|x|}
 \le \frac4{\sqrt n}.                                  \tag{1}
\end{equation}
In particular,
$d_e(f_n,f_m)\le4(n^{-1/2}+m^{-1/2})$, so $(f_n)$ is $d_e$-Cauchy.

The limit $f$ is continuous, compactly supported, and extensively bounded:
on its support $|x|\ge\tfrac14$ and $|f(x)|\le\tfrac12$, hence
$|f(x)|/|x|\le2$.  It is not Lipschitz.  Indeed, for
$0<h<\tfrac14$,
\[
 f(a)=0,\qquad f(a+h)=(1-4h)\sqrt h,\qquad
 \frac{|f(a+h)-f(a)|}{h}=\frac{1-4h}{\sqrt h}\longrightarrow\infty.
\]
Therefore the $d_e$-limit in the complete metric space $E(\R,\R)$ does not
belong to $\Lip_0(\R,\R)$.  Uniqueness of limits rules out any other
$\Lip_0$ limit, proving Theorem~\ref{thm:main}. \qed

\begin{corollary}
For every nonzero Banach space $Y$, the space $\Lip_0(\R,Y)$ is incomplete
under $d_e$.
\end{corollary}
\begin{proof}
Fix $y\in Y$ with $\|y\|=1$ and replace $f_n,f$ by $f_ny,fy$.
All estimates and the cusp obstruction are unchanged.
\end{proof}

\section*{Proof Intuition}

The metric $d_e$ controls the size of a map only relative to the single
distinguished point.  On a region bounded away from zero, it is merely a
weighted uniform metric.  Uniform convergence preserves continuity but not
a uniform Lipschitz constant.  We therefore regularize a square-root cusp
inside such a region.  The Lipschitz constants of the approximants diverge,
but $d_e$ sees only their uniformly vanishing vertical error.

This also shows that the failure occurs inside the source paper's complete
continuous subspace $E_c(\R,\R)$: continuity survives, while global
Lipschitz regularity does not.

\section*{Scope, novelty audit, and review recommendation}

The theorem fully answers the universal question negatively, with both
domain and codomain Banach spaces.  It does not characterize the pointed
metric spaces for which completeness does hold.

Exact-question, exact-title, author, arXiv-id, $\Lip_0(M,N)$, $d_e$, and
``extensively bounded'' searches through 13 August 2026 found the source and
a current index still listing the question as open, but no later explicit
answer or equivalent counterexample.  This is a bounded novelty audit, not a
priority claim.  Human review should check the definition match, estimate
\eqref{conv}, and the cusp quotient; no external theorem is used.

\section*{Reference}

\textbf{[1]} A. K. Karn and A. Mandal,
\emph{A generalization of Lipschitz mappings}, arXiv:2410.10677,
Remark 3.3, PDF p.~11.

\end{document}
